\section{Fourier credit and the outcome plateau}
\label{app:outcome}
\subsection{Fourier contraction}
\label{app:C-fourier}
\label{app:assumptions}




(A1)--(A3) denote \cref{ass:arch,ass:init,ass:task}, and (A3$'$) a uniform start. Unless stated otherwise the query and the start
count are fixed. Constants $c,C$ depend only on $a_0,\embconst,K$ and the numbers of blocks and heads. Every bound stated on $\initball$ also
holds on $\initball'=\{\bm\theta:\|\bm\theta-\bm\theta_0\|\le2\ballradius\initscale,\ \beta=0\}$, whose relative radius is at most $\min(\frac{1}{4},\frac{\embconst}{2})$; there
$|b_m|\ge(1-2\ballradius)\initscale\ge\frac{3}{4}\initscale$ and $\sigma_{\min}(\mathbf E_{\mathrm{val}})\ge\frac{\initscale\embconst}{2}$. The only width-dependent threshold is $\initscale_0=c\min(1,\frac{1}{N_b})$.

\begin{lemma}[Gradient split]
\label{lem:gradient-split}
For every differentiable model, with $\outcredit(\bm\theta)=\E[\langle\bm z_{\bm\theta},\bm y-\bm u\rangle]$,
\[
 \nabla_{\bm\theta}\Lout=\underbrace{\E\bigl[(\Jz)^\top(\bm p-\bm u)\bigr]}_{\text{answer-independent}}-\underbrace{\nabla_{\bm\theta}\outcredit}_{\text{answer-dependent}},\qquad
 \nabla_{\bm\theta}\outcredit=\E\bigl[(\Jz)^\top(\bm y-\bm u)\bigr].
\]
\end{lemma}

\begin{proof}
The logit gradient of the cross-entropy is $\bm p-\bm e_y=(\bm p-\bm u)-(\bm y-\bm u)$.
\end{proof}

\paragraph{Conditioning on unseen letters.}
\label{app:target-spectrum}

\paragraph{Contraction.}\label{app:contraction} By \cref{eq:lambda}, $\mixrate^2=1-2p(1-p)(1-\cos\tfrac{2\pi}{K})$, and for $K=2$,
$\mixrate=|1-2p|$ and $1-\mixrate^2=4p(1-p)$. With $\alpha=e^{2\pi\mathrm i/K}$ and the Fourier multipliers
$\rho_j=1-p+p\alpha^j$, so that $\mixrate=\max_{j\ne0}|\rho_j|$, one unseen letter multiplies each nonconstant count
mode $\alpha^{js_t}$ by $\rho_j$, while one displayed transition keeps the innovation
\begin{equation}
 \E\bigl[\bigl|\alpha^{j s_t}-\rho_j\,\alpha^{j s_{t-1}}\bigr|^2\bigr]=1-|\rho_j|^2\ge1-\mixrate^2 .
 \label{eq:move-variance}
\end{equation}
In the Fourier basis of $\one^\perp$, $\operatorname{Cov}(\bm g_t)$ is diagonal with entries $\tfrac{1}{K}(1-|\rho_j|^2)=\tfrac{2p(1-p)}{K}\bigl(1-\cos\tfrac{2\pi j}{K}\bigr)$,
which bounds the covariance below by $(1-\lambda^2)/K$ on $\mathbf1^\perp$.

\begin{lemma}[Conditional contraction]
\label{lem:target-contraction}
Assume (A3) and fix the start count and the query, so that the letters are independent. Let $\bm\xi=\bm y-\bm u$, $S\subseteq[n]$ and $\mathcal F_S=\sigma(w_t:t\in S)$. Then
\begin{equation}
 \bigl\|\E[\bm\xi\mid\mathcal F_S]\bigr\|^2=\frac{1}{K}\sum_{j=1}^{K-1}|\rho_j|^{2(n-|S|)}\le\Bigl(1-\frac{1}{K}\Bigr)\mixrate^{2(n-|S|)} .
 \label{eq:horizon}
\end{equation}
Hence for every square-integrable $\mathcal F_S$-measurable matrix $\mathbf H$ with $K$ rows,
 $\bigl\|\E[\mathbf H^\top\bm\xi]\bigr\|\le\mixrate^{n-|S|}\,\|\mathbf H\|_{L^2},\qquad \|\mathbf H\|_{L^2}:=\bigl(\E\|\mathbf H\|_F^2\bigr)^{1/2}.$ 
\end{lemma}

\begin{proof}
Write $s_n=s_0+\sum_tb_t$. The letters are independent, so for $j\ne0$
\[
 \E\bigl[\alpha^{js_n}\bigm|\mathcal F_S\bigr]=\alpha^{js_0}\prod_{t\in S}\alpha^{jb_t}\prod_{t\notin S}\E\bigl[\alpha^{jb_t}\bigr]=\alpha^{js_0}\,\rho_j^{\,n-|S|}\prod_{t\in S}\alpha^{jb_t},
\]
whose modulus is $|\rho_j|^{n-|S|}$. Parseval over the nonconstant modes $j=1,\ldots,K-1$ gives the identity, and $|\rho_j|\le\mixrate$ the bound.
The matrix bound follows from \cref{eq:horizon} by the tower property
and Cauchy--Schwarz.
\end{proof}

\paragraph{Sharpness.}\label{rem:tight} The $j=1$ term of \cref{eq:horizon} is $\frac{1}{K}\mixrate^{2(n-|S|)}$, and $\mathbf H=\E[\bm\xi\mid\mathcal F_S]$ attains the second bound up to the
factor $(1-\frac{1}{K})^{1/2}$, so the exponent $n-|S|$ is sharp for $0<\mixrate<1$.

\subsection{Transformer score expansion}
\label{app:model-spectrum}

The contraction estimate helps only if we can control how many letters a network term uses. Expanding the attention scores does
this without asking the entire output to depend on a small set: each term has its own set of letters, and terms with larger sets
are smaller near initialization.

\paragraph{Score expansion.} For $\zeta\in[0,1]$, let $\bm z(\zeta)$ be the network with every score $\psi_{il}$, in both blocks, replaced by
$\zeta\psi_{il}$. Then $\bm z(1)=\bm z_{\bm\theta}$, and for every integer $k_0\ge0$,
\begin{equation}
 \bm z_{\bm\theta}=\sum_{k=0}^{k_0}\bm z^{[k]}+\bm z_{\mathrm{rem}},\qquad
 \bm z^{[k]}=\frac{1}{k!}\frac{d^k\bm z}{d\zeta^k}(0),\qquad
 \bm z_{\mathrm{rem}}=\int_0^1\frac{(1-\zeta)^{k_0}}{k_0!}\frac{d^{k_0+1}\bm z}{d\zeta^{k_0+1}}(\zeta)\,d\zeta .
 \label{eq:score-split}
\end{equation}
$\mathbf J^{[k]}$ and $\mathbf J_{\mathrm{rem}}$ are the corresponding Jacobian terms. A term \emph{reads} the letters whose inputs it contains.

\begin{proposition}[Score count and letter support]
\label{prop:score-degree}
Fix the start state and the query. For $\bm\theta\in\initball$, $\Jz=\sum_{k=0}^{k_0}\mathbf J^{[k]}+\mathbf J_{\mathrm{rem}}$, and there is $\attconst$ such that:
\begin{enumerate}[label=(\roman*),leftmargin=*,itemsep=2pt,topsep=2pt]
\item $\mathbf J^{[0]}=\mathbf J^{[0]}_\emptyset+\sum_{t=1}^n\mathbf J^{[0]}_t$, where $\mathbf J^{[0]}_\emptyset$ reads no letter, $\mathbf J^{[0]}_t$ reads only $w_t$, $\sup\|\mathbf J^{[0]}_\emptyset\|_F\le C\initscale$ and
$\sum_t\sup\|\mathbf J^{[0]}_t\|_F\le C\initscale^2$;
\item for $k\ge1$, $\mathbf J^{[k]}=\sum_\nu\mathbf T_{k,\nu}$, where each $\mathbf T_{k,\nu}$ reads at most $2k+1$ letters and $\sum_\nu\sup\|\mathbf T_{k,\nu}\|_F\le C\initscale^2(\attconst\initscale^4)^k$;
\item $\sup\|\mathbf J_{\mathrm{rem}}\|_F\le C\initscale^2(\attconst\initscale^4)^{k_0+1}$.
\end{enumerate}
The same bounds hold for $\bm z^{[k]}$ and $\bm z_{\mathrm{rem}}$ with one more factor $\initscale$. Suprema are over inputs; constants do not depend on
$n$, $k$ or $k_0$.
\end{proposition}

$\attconst$ is fixed once, large enough for \cref{prop:score-degree,prop:network}; it depends only on $a_0$, $\embconst$ and $K$.

\label{app:network}

\begin{proposition}[Bounds for the score expansion]
\label{prop:network}
Assume (A1)--(A2). Uniformly on $\initball$, at every position and input, with $\bm x_i$ the token-plus-position vector,
all derivatives with respect to $s$ are evaluated at $s=0$ along $\bm\theta+sv$.
Let $\|F\|_{\mathrm{abs}}$ be the absolute sum of $F$ at $\zeta=0$: expand $F$ into monomials, products of
parameter blocks, inputs $\bm x_j$ and input-independent attention weights; bound each monomial by the product of the norm bounds of its
parameter factors on the ball, of $\sup\|\bm x_j\|$ for its inputs, and of its weights; and sum these bounds before any cancellation:
\begin{enumerate}[label=(\roman*),leftmargin=*,itemsep=3pt,topsep=4pt]
\item with input-independent offsets $\bm O^1=\bm o_1$, $\bm O^2=\bm o_1+\bm o_2$: $\bm h_i^\ell=\bm x_i+\bm O^\ell+O(\initscale^2)$, $|\psi_{il}|\le C\initscale^4$,
$\|\bm z_i\|\le C\initscale^2$, $\|\dd{\bm z_i}{s}\|\le C\initscale\|v\|$;
\item for $k\ge1$, $\zeta\in[0,1]$ and $\|v\|=1$,
\[
 \frac{1}{k!}\Bigl\|\frac{d^k\bm z_i}{d\zeta^k}\Bigr\|\le C\initscale^3(C_0\initscale^4)^k,\qquad
 \frac{1}{k!}\Bigl\|\dd{}{s}\left(\frac{\mathrm d^k\bm z_i}{\mathrm d\zeta^k}\right)\Bigr\|\le C\initscale^2(C_0\initscale^4)^k;
\]
the same holds for the absolute sums at $\zeta=0$, for $\bm h^2_i$ with one power of $\initscale$ fewer, and for the Taylor remainder with $k=k_0+1$;
\item the zero-score network is affine in the inputs; its transported logit terms have tangent $O(\initscale^2)$, and direct-input and
input-free terms have tangent $O(\initscale)$; an order-$k$ coefficient contains at most $2k+1$ inputs; a one-score logit term has size
$O(\initscale^7)$ and tangent $O(\initscale^6)$;
\end{enumerate}
The constants do not depend on $n$, the widths or $k$.
\end{proposition}

\begin{proof}
For a network quantity $F$ let $F^{[k]}(\zeta)=\frac{1}{k!}\frac{d^kF}{d\zeta^k}(\zeta)$, so that $\bm z^{[k]}=\bm z^{[k]}(0)$, and write $\|F^{[k]}\|$ for its supremum over
inputs, positions and $\zeta\in[0,1]$ (the $\ell_1$ norm for attention rows).

\emph{(i).} $|b_m|\ge\frac{3}{4}\initscale>C\initscale^2\ge|(\mathbf W_{\mathrm{in}}\bm h)_m|$, so one activation matrix $\mathbf D_{\mathrm{act}}$ serves every input and every $\zeta$, and
\[
 \begin{aligned}
 \bm h_i^\ell&=\mathbf M_\ell\bigl(\bm h_i^{\ell-1}+\bm a_i^\ell\bigr)+\bm o_\ell,\\
 \mathbf M_\ell&=\mathbf I+\mathbf W^\ell_{\mathrm{out}}\mathbf D_{\mathrm{act}}\mathbf W^\ell_{\mathrm{in}},\qquad
 \bm o_\ell=\mathbf W^\ell_{\mathrm{out}}\mathbf D_{\mathrm{act}}\bm b^\ell.
 \end{aligned}
\]
with $\bm a_i^\ell=\sum_{\mathrm{heads}}\sum_{j\le i}\omega_{ij}\mathbf W_V\bm h^{\ell-1}_j$ and $\omega_{i\cdot}=\softmax(\zeta\psi_{i\cdot})$. Then
 \[
\begin{gathered}
\|\mathbf M_\ell-\mathbf I\|\le C\initscale^2,\\ \|\bm o_\ell\|\le CN_b\initscale^2\le C\initscale,\\ \|\dd{\bm o_\ell}{s}\|\le(a_0+2N_b)\initscale\le C,
\end{gathered}
\] 
and $\|\omega_{i\cdot}\|_1=1$. Induction over the two blocks gives (i).

\emph{Softmax calculus.} Let $\omega=\softmax(s)$ and
$\E_\omega[f]=\sum_j\omega_jf_j$. For a direction $\bm\nu$,
\begin{align*}
 \left.\dd{}{t}\omega_j(s+t\bm\nu)\right|_{t=0}
 &=\omega_j(\nu_j-\E_\omega[\bm\nu]),\\
 \left.\dd{}{t}\E_{\omega(s+t\bm\nu)}[f]\right|_{t=0}
 &=\E_\omega[\bm\nu f]-\E_\omega[\bm\nu]\E_\omega[f].
\end{align*}
After $k$ derivatives, each term has at most $k+1$ differentiable factors
and contains each direction once. The absolute coefficient sums satisfy
$a_{k+1}\le2(k+1)a_k$, $a_0=1$. Summing against the probability weights gives
\begin{equation}
 \bigl\|D^k\omega(s)[\bm\nu_1,\ldots,\bm\nu_k]\bigr\|_1\le2^kk!\prod_a\|\bm\nu_a\|_\infty .
 \label{eq:softmax-avg}
\end{equation}
For the smooth finite-dimensional maps in the fixed-ReLU region, the
higher-order chain rule \citep[Theorem~2]{clark2013faa} yields
\[
 \frac{1}{j!}\|D^jF\|\le Ac^j,\quad \|h^{[k]}\|\le ab^k
 \quad\Longrightarrow\quad
 \|(F\circ h)^{[k]}\|\le Aca\,b^k(1+ca)^{k-1},\quad k\ge1.
\]
Here the inner map's derivative bounds hold for every positive order; the
composition count is $\binom{k-1}{j-1}$.

\emph{(ii), values.} With $\|\bm\omega^{\ell,[0]}\|_1=1$, $\|\bm h^{0,[k]}\|=0$ for $k\ge1$, and $a_r\ge1$ in the compositions,
\begin{alignat*}{2}
 \|\psi^{\ell,[k]}\|&\le C\initscale^2\textstyle\sum_{a+b=k}\|\bm h^{\ell-1,[a]}\|\,\|\bm h^{\ell-1,[b]}\| &\qquad& \why{score is a product}\\
 \|(\zeta\psi^\ell)^{[k]}\|&\le\|\psi^{\ell,[k]}\|+\|\psi^{\ell,[k-1]}\| && \why{$(\zeta\psi)^{[k]}=\zeta\psi^{[k]}+\psi^{[k-1]}$}\\
 \|\bm\omega^{\ell,[k]}\|_1&\le\textstyle\sum_{j=1}^k2^j\sum_{a_1+\cdots+a_j=k}\prod_r\|(\zeta\psi^\ell)^{[a_r]}\| && \why{composition, \cref{eq:softmax-avg}}\\
 \|\bm a^{\ell,[k]}\|&\le C\initscale\textstyle\sum_{a+c=k}\|\bm\omega^{\ell,[a]}\|_1\,\|\bm h^{\ell-1,[c]}\| && \why{head output is a product}\\
 \|\bm h^{\ell,[k]}\|&\le C\bigl(\|\bm h^{\ell-1,[k]}\|+\|\bm a^{\ell,[k]}\|\bigr) && \why{affine block}\\
 \|\bm z^{[k]}\|&\le C\initscale\,\|\bm h^{2,[k]}\| && \why{readout}
\end{alignat*}
With $b=C\initscale^4$, block $1$ gives, for $k\ge1$, $\|\bm\omega^{1,[k]}\|_1\le b^k$ and $\|\bm a^{1,[k]}\|+\|\bm h^{1,[k]}\|\le C\initscale^2b^k$ (while $\|\bm h^{1,[0]}\|\le C\initscale$). Then $\|\psi^{2,[0]}\|\le C\initscale^4$,
$\|\psi^{2,[k]}\|\le C\initscale^5(k+1)b^k$ for $k\ge1$, and $\|(\zeta\psi^2)^{[k]}\|\le C(2b)^k$, so
\begin{gather*}
 \|\bm\omega^{2,[k]}\|_1\le C\sum_{j=1}^k\binom{k-1}{j-1}C^j(2b)^k\le C(Cb)^k,\\
 \|\bm a^{2,[k]}\|+\|\bm h^{2,[k]}\|\le C\initscale^2(Cb)^k,\qquad \|\bm z^{[k]}\|\le C\initscale^3(C_0\initscale^4)^k .
\end{gather*}

\emph{(ii), tangents, remainders, absolute sums.} Differentiating along $v$ replaces one factor of each product by its tangent, or one
$(\zeta\psi)^{[a]}$ in a softmax term, or adds $D^{j+1}\omega[\dd{(\zeta\psi)^{[0]}}{s},\ldots]$. With $\|\dd{\bm h^{0,[0]}}{s}\|=O(1)$, $\|\dd{\psi^{\ell,[0]}}{s}\|=O(\initscale^3)$
and $j+1\le2^k$ absorbed into $C_0^k$,
\begin{gather*}
 \|\dd{\bm\omega^{\ell,[k]}}{s}\|_1\le C\initscale^{-1}(C_0\initscale^4)^k,\qquad \|\dd{\bm a^{\ell,[k]}}{s}\|+\|\dd{\bm h^{\ell,[k]}}{s}\|\le C\initscale(C_0\initscale^4)^k,\\
 \|\dd{\bm z^{[k]}}{s}\|\le C\initscale^2(C_0\initscale^4)^k .
\end{gather*}
The integral remainder is bounded by the supremum of the next normalized derivative.

\emph{Absolute sums.} $\|\cdot\|_{\mathrm{abs}}$ is subadditive and submultiplicative, and $\|F\|\le\|F\|_{\mathrm{abs}}$. At $\zeta=0$ every attention row is an
input-independent probability vector $\bm\pi$ (uniform over the visible sites here, $\softmax(\beta)$ in \cref{lem:bias-uniform}), and
 $\Bigl\|\textstyle\sum_j\pi_j\prod_rF_r(j)\Bigr\|_{\mathrm{abs}}\le\prod_r\max_j\|F_r(j)\|_{\mathrm{abs}} .$ 
At $\zeta=0$, $D^j\omega$ expands as $\omega(0)$-weights times products of direction entries and $\omega(0)$-averages, with total coefficient at most $2^jj!$,
so \cref{eq:softmax-avg} holds for absolute sums with $\|\bm\nu_a\|_\infty$ replaced by $\max_l\|\nu_{a,l}\|_{\mathrm{abs}}$. Hence every line of the recursion
above holds verbatim for $\|\cdot\|_{\mathrm{abs}}$ at $\zeta=0$, which gives the absolute-sum bounds of (ii).

\emph{Tangents of absolute sums.} Take the norm bound of every parameter factor at least $\initscale$ (weights and embeddings $(a_0+1)\initscale$,
biases $(N_b+1)\initscale$). An order-$k$ monomial has at most $C(k+1)$ factors: at most $k$ scores, each with $\mathbf W_Q$ and $\mathbf W_K$, and at
most $2k+1$ inputs, each carried by at most four parameter blocks. Replacing one factor by its unit-direction tangent multiplies the bound by at
most $\frac{1}{\initscale}$, so the tangent bound is at most $\frac{C(k+1)}{\initscale}\le\frac{C2^k}{\initscale}$ times the monomial bound, uniformly in the unit direction
$v$; this lets $\sup_v$ pass inside the sum over monomials.

\emph{(iii).} At zero scores the network is affine in the inputs (for a head with a submask, $\frac{1}{m_j}$ becomes one over the size of its
visible set):
\begin{gather*}
 \bm h^1_j=\mathbf M_1\Bigl(\bm x_j+\frac{1}{m_j}\sum_{l\le j}\mathbf W^1_V\bm x_l\Bigr)+\bm o_1,\\
 \bm z^{[0]}_i=\mathbf V\Bigl[\mathbf M_2\Bigl(\bm h^1_i+\frac{1}{m_i}\sum_{j\le i}\mathbf W^2_V\bm h^1_j\Bigr)+\bm o_2\Bigr].
\end{gather*}
Transported terms contain $\mathbf V$, $\mathbf W_V$ and some $\bm x_j$, so their tangent is $O(\initscale^2)$; direct-input and input-free terms have tangent $O(\initscale)$.
With $d(F)$ the number of input occurrences in $F$,
\begin{alignat*}{2}
 d(\bm h^{\ell,[a]})&\le2a+1 &\qquad& \why{induction on $\ell$}\\
 d(\psi^{\ell,[b]})&\le2b+2 && \why{score is a product of two residuals}\\
 d\bigl((\zeta\psi)^{[a]}\bigr)&\le2a && \why{$(\zeta\psi)^{[a]}=\psi^{[a-1]}$}\\
 d(\text{order-$k$ head term})&\le2a_1+\cdots+2a_j+2c+1=2k+1 && \why{$a_1+\cdots+a_j+c=k$}
\end{alignat*}
and affine MLPs and parameter derivatives add none. The one-score correction is
 $\frac{1}{m_i}\sum_{j\le i}\bigl(\psi_{ij}(0)-\bar\psi_i(0)\bigr)\mathbf W_V\bm h_j(0),\qquad \bar\psi_i(0)=\frac{1}{m_i}\sum_{l\le i}\psi_{il}(0);$ 
its logit terms contain $\mathbf V,\mathbf W_V,\mathbf W_Q,\mathbf W_K$ and three residual inputs, so they have size $O(\initscale^7)$ and tangent $O(\initscale^6)$.

\end{proof}

\begin{proof}[Proof of \cref{prop:score-degree}]
Expand each coefficient into monomials in parameters and inputs. A Jacobian entry of a monomial is a monomial with the same inputs, and the
Frobenius norm of a map into $\R^K$ is at most $\sqrt K$ times its operator norm.

\noindent(i) At zero scores the network is affine in the inputs (\cref{prop:network}(iii)): each term contains at most one input. Group the terms by
that input. At the answer position the direct input is fixed and contains no letter; the transported terms have tangent $O(\initscale^2)$, and
their coefficients sum absolutely to a constant because every zero-score attention row is a probability vector.

\noindent(ii) By \cref{prop:network}(iii) an order-$k$ monomial contains at most $2k+1$ inputs, and by \cref{prop:network}(ii) the absolute sum of the
tangents at $\zeta=0$ is at most $C\initscale^2(C_0\initscale^4)^k$; take $\attconst\ge C_0$.

\noindent(iii) Apply \cref{prop:network}(ii) with $k=k_0+1$. Logit bounds carry one more factor $\initscale$ than Jacobian bounds.
\end{proof}

\subsection{Exponential outcome plateau}
\label{app:thm1}

Combining \cref{lem:target-contraction,prop:network} gives the following credit bound.

\begin{lemma}[Credit summation]
\label{lem:score-sum}
Assume (A1)--(A3), $\bm\theta\in\initball$, $0<\initscale\le\initscale_0$ and $0<\mixrate<1$, with the start state and query fixed. Let
$k_0=\lceil\frac{n-1}{2}\rceil$. If $\attconst\initscale^4\le\frac{1}{2}\mixrate^2$, then
 $\bigl\|\E\bigl[\mathbf J_{\mathrm{score}}^\top(\bm y-\bm u)\bigr]\bigr\|\le C\initscale^6\mixrate^{n-3},\qquad \mathbf J_{\mathrm{score}}=\Jz-\mathbf J^{[0]}.$ 
If moreover $\initscale\le\mixrate$, then
\begin{equation}
 \|\nabla_{\bm\theta}\outcredit(\bm\theta)\|\le C\initscale\mixrate^n,\qquad |\outcredit(\bm\theta)|\le C\initscale^2\mixrate^n.
 \label{eq:thm1-full}
\end{equation}
The bounds persist after averaging over the start state and query.
\end{lemma}

\begin{proof}
Apply \cref{lem:target-contraction} to each term of \cref{prop:score-degree}, with $\|\mathbf H\|_{L^2}\le\sup\|\mathbf H\|_F$ and $\|\bm y-\bm u\|\le1$ for the remainder:
\begin{alignat*}{2}
 \bigl\|\E\bigl[(\mathbf J^{[0]})^\top(\bm y-\bm u)\bigr]\bigr\|&\le C\initscale\mixrate^n+C\initscale^2\mixrate^{n-1} &\qquad& \why{(i): $0$ or $1$ letter}\\
 \bigl\|\E\bigl[(\mathbf J^{[k]})^\top(\bm y-\bm u)\bigr]\bigr\|&\le C\initscale^2(\attconst\initscale^4)^k\mixrate^{n-2k-1} && \why{(ii): $\le2k+1$ letters, $k\ge1$}\\
 &=C\initscale^6\mixrate^{n-3}\Bigl(\frac{\attconst\initscale^4}{\mixrate^2}\Bigr)^{k-1} && \\
 \bigl\|\E\bigl[\mathbf J_{\mathrm{rem}}^\top(\bm y-\bm u)\bigr]\bigr\|&\le C\initscale^2(\attconst\initscale^4)^{k_0+1} && \why{(iii)}\\
 &=C\attconst\initscale^6\mixrate^{2k_0}\Bigl(\frac{\attconst\initscale^4}{\mixrate^2}\Bigr)^{k_0} && \\
 &\le C\initscale^6\mixrate^{n-3}\Bigl(\frac{\attconst\initscale^4}{\mixrate^2}\Bigr)^{k_0} && \why{$2k_0\ge n-1$}
\end{alignat*}
If $2k+1>n$, the factor $\mixrate^{n-2k-1}\ge1$ still bounds the credit. Since $\frac{\attconst\initscale^4}{\mixrate^2}\le\frac{1}{2}$, the
score-dependent terms sum to at most $3C\initscale^6\mixrate^{n-3}$, which is the first bound. Adding the zero-order term,
\begin{align*}
 \|\nabla_{\bm\theta}\outcredit\|&\le C\bigl(\initscale\mixrate^n+\initscale^2\mixrate^{n-1}+\initscale^6\mixrate^{n-3}\bigr),\\
 |\outcredit|&\le C\bigl(\initscale^2\mixrate^n+\initscale^3\mixrate^{n-1}+\initscale^7\mixrate^{n-3}\bigr),
\end{align*}
by \cref{lem:gradient-split} and the same argument for the logits, and $\initscale\le\mixrate$ reduces both to \cref{eq:thm1-full}.
\end{proof}

For $c$ small enough, $\initscale\le c\mixrate$ implies $\initscale\le\mixrate$ and $\attconst\initscale^4\le\frac{1}{2}\mixrate^2$, and $\initscale^2\le c\mixrate$ implies $\attconst\initscale^4\le\frac{1}{2}\mixrate^2$.



\paragraph{A trained relative bias.} Let $\bm\vartheta$ denote the parameters other than $\beta$, and
$\mathcal B^{\mathrm{full}}_{\initscale}=\{\bm\theta:\|\bm\vartheta-\bm\vartheta_0\|\le2\ballradius\initscale\}$ with $\beta$ unrestricted. A zero readout in $\tilde{\bm\theta}$ is admissible in (A2).

\begin{lemma}[Uniformity in the relative bias]
\label{lem:bias-uniform}
Assume (A1)--(A3), $\initscale\le\initscale_0$. On $\mathcal B^{\mathrm{full}}_{\initscale}$, with constants independent of $\beta$:
\begin{enumerate}[label=(\alph*),leftmargin=*,itemsep=3pt,topsep=4pt]
\item the logit bounds of \cref{prop:network} and of \cref{prop:score-degree} for $\bm z^{[k]},\bm z_{\mathrm{rem}}$ hold;
\item if $\initscale\le c\mixrate$, then $|\outcredit(\bm\theta)|\le C_2\initscale^2\mixrate^n$;
\item the conclusions of \cref{lem:end-smooth} hold for unit directions in all parameters, $\beta$ included.
\end{enumerate}
\end{lemma}

\begin{proof}
Each attention row is
 $\bm\omega_{i\cdot}(\zeta)=\softmax\bigl(\beta_{i-\cdot}+\zeta\psi_{i\cdot}\bigr)\qquad\text{over its visible sites.}$ 

\noindent(a) At $\zeta=0$,
 $\bm\omega_{i\cdot}(0)=\softmax(\beta_{i-\cdot}),\qquad \bm\omega_{i\cdot}(0)\ge0,\qquad \|\bm\omega_{i\cdot}(0)\|_1=1,$ 
which does not depend on the input; the zero-score network is affine with the same input counts $d(\cdot)$. Since $\beta$ does not
depend on $\zeta$,
 $(\beta+\zeta\psi)^{[k]}=\zeta\psi^{[k]}+\psi^{[k-1]}\qquad(k\ge1),$ 
so the recursion of \cref{prop:network}(ii) is unchanged, and \cref{eq:softmax-avg} holds uniformly in $\beta$. The ReLU pattern stays
fixed, since $|(\mathbf W_{\mathrm{in}}\bm h)_m|\le C\initscale^2<|b_m|$ for any attention weights. The uniform weights $\frac{1}{m_j}$ enter only
\cref{prop:network}, which is not needed here.

\smallskip\noindent(b) The proof of \cref{lem:score-sum} uses only \cref{lem:target-contraction} and the logit bounds of (a), so
$|\outcredit(\bm\theta)|\le C_2\initscale^2\mixrate^n$.

\smallskip\noindent(c) A unit direction with $\beta$-component $v_\beta$ perturbs each row's scores by $\bm\nu=v_\beta$, $\|\bm\nu\|_\infty\le1$. By
\cref{eq:softmax-avg} with $k=1,2$, the first two row derivatives have $\ell_1$ norms at most $C$. The product-rule logit bounds in \cref{lem:end-smooth} therefore apply unchanged.
\end{proof}

\begin{lemma}[Smoothness near initialization]
\label{lem:end-smooth}
Assume (A1)--(A2), $\initscale\le\initscale_0$. There is $C_1$, independent of $n$, $\initscale$ and the widths, such that for every $\bm\theta\in\initball'$ and
every unit direction $\bm v$,
 $\|\nabla_{\bm\theta}\Lout(\bm\theta)\|\le C_1\initscale,\qquad \Bigl|\frac{d^2}{ds^2}\Lout(\bm\theta+s\bm v)\big|_{s=0}\Bigr|\le C_1 .$ 
\end{lemma}

\begin{proof}
Along $\bm\theta+s\bm v$ at $s=0$, applying the product rule twice to the
two blocks and \cref{eq:softmax-avg} with $k=1,2$ gives
\[
 \|\bm z\|\le C\initscale^2,\qquad
 \left\|\dd{\bm z}{s}\right\|\le C\initscale,\qquad
 \left\|\frac{\mathrm d^2\bm z}{\mathrm d s^2}\right\|\le C.
\]
The ReLU pattern is fixed by \cref{prop:network}(i). These bounds use only
operator norms and attention-row $\ell_1$ norms bounded by one, so $C$ is
independent of $n$ and the widths.
For cross-entropy, the log-sum-exp derivative formula
\citep[Section~3.1.5]{boyd2004convex} gives
\[
 \nabla_{\bm z}\ell=\bm p-\bm e_y,\qquad
 \nabla_{\bm z}^2\ell=\operatorname{diag}\bm p-\bm p\bm p^\top,
 \qquad \|\nabla_{\bm z}\ell\|\le\sqrt2,\quad
 \|\nabla_{\bm z}^2\ell\|\le1.
\]
The chain rule therefore yields
$|\dd{\ell}{s}|\le C\initscale$ and
$|\mathrm d^2\ell/\mathrm d s^2|\le C$.
Average over prompts and take the supremum over unit $\bm v$.
\end{proof}

\paragraph{Outcome plateau.}\label{app:end} For one prompt, $\KL(\bm u\|\bm p)=\log\sum_ae^{z_a}-\log K-\langle\bm u,\bm z\rangle$, hence
\begin{equation}
 \Lout(\bm\theta)=\log K+\E\bigl[\KL(\bm u\,\|\,\bm p_{\bm\theta})\bigr]-\outcredit(\bm\theta).
 \label{eq:lo-identity}
\end{equation}

\paragraph{Ambient projected flow.}
For this supporting estimate, let $P(t)$ be an orthogonal projection onto trainable ambient parameter blocks, with finitely many switches on each bounded interval, and let $\dot\theta=-P(t)\nabla_\theta\Lout$ start at the zero-readout seed. The calculation below gives the local loss, displacement and assigned-probability bounds until time $r_0^2\lambda^{-n}/C_2$. Full ambient joint flow is included. Appendix~\ref{app:match} supplies the separate metric comparison required for the tied scalar flow in Theorem~\ref{thm:matched}.

\begin{proof}[Ambient outcome plateau]
\emph{Loss floor.} Let $B_n=C_2\initscale^2\mixrate^n$, $\initscale\le c\mixrate$ and $\bm\theta\in\mathcal B^{\mathrm{full}}_{\initscale}$.
\begin{alignat}{2}
 \Lout(\bm\theta)&=\log K+\E\,\KL(\bm u\,\|\,\bm p_{\bm\theta})-\outcredit(\bm\theta) &\qquad& \why{\cref{eq:lo-identity}}\nonumber\\
 &\ge\log K-B_n && \why{$\KL\ge0$, \cref{lem:bias-uniform}(b)}\label{eq:lo-floor}
\end{alignat}

\emph{Escape time.} By \cref{lem:end-smooth} and \cref{lem:bias-uniform}(c), $\Lout$ is smooth on $\mathcal B^{\mathrm{full}}_{\initscale}$, so the flow is defined between the
switching times of $\mathbf P(t)$ and continuous across them. Let $T_{\mathrm{exit}}=\inf\{t:\|\bm\theta(t)-\bm\theta_0\|>\ballradius\initscale\}$, the norm
including $\beta$. For $t\le T_{\mathrm{exit}}$,
\begin{alignat*}{2}
 \ddt{\Lout}&=-\|\mathbf P\nabla_{\bm\theta}\Lout\|^2=-\Bigl\|\ddt{\bm\theta}\Bigr\|^2 &\qquad& \why{$\mathbf P$ constant between switches}\\
 \int_0^t\Bigl\|\ddt{\bm\theta}\Bigr\|^2&=\Lout(\bm\theta_0)-\Lout(\bm\theta(t))\le B_n && \why{$\Lout(\bm\theta_0)=\log K$, \cref{eq:lo-floor}}\\
 \|\bm\theta(t)-\bm\theta_0\|^2&\le t\int_0^t\Bigl\|\ddt{\bm\theta}\Bigr\|^2\le t\,B_n && \why{Cauchy--Schwarz}\\
 T_{\mathrm{exit}}&\ge\frac{\ballradius^2}{C_2}\,\mixrate^{-n} && \why{$\|\bm\theta(T_{\mathrm{exit}})-\bm\theta_0\|=\ballradius\initscale$}
\end{alignat*}
and \cref{eq:lo-floor} holds on $[0,T_{\mathrm{exit}}]$.

\emph{Near-chance prediction.}
\begin{alignat*}{2}
 \E\,\KL(\bm u\,\|\,\bm p_t)&\le\outcredit(\bm\theta(t))\le B_n &\qquad& \why{\cref{eq:lo-identity}, $\Lout(\bm\theta(t))\le\log K$}\\
 \E[\bm p_{t,y}]-\tfrac{1}{K}&\le\E\,\mathrm{TV}(\bm p_t,\bm u)\le\sqrt{\frac{C_2}{2}}\;\initscale\mixrate^{n/2} && \why{Pinsker, Jensen}
\end{alignat*}

\emph{Constants.} Take $c$ the smaller of $\frac{\ballradius^2}{C_2}$ and the constant for which $\initscale\le c\mixrate$ gives $\initscale\le\mixrate$ and
$\attconst\initscale^4\le\frac{1}{2}\mixrate^2$, and $C=\max(C_2,\sqrt{\frac{C_2}{2}})$.
\end{proof}
